\section{When It Fails Instead}
\label{sec:ch18-when-it-fails-instead}

\index{blame routing!a failure, routed|(}%
Section~\ref{sec:ch18-writing-it-down} found a slow service by counting what
it did. A failing service is the easier half of the problem, because the
router already names it, and the harder half of the fix, because naming it is
where the router stops.

\index{serviceName@\code{serviceName}!the router names the failing service}%
Take a failure that needs no bug in the service's code: the Ratings service
running against a database whose \code{Ratings} table is not there, which is
the state a missed migration leaves behind. Rename the table out from under the running
service, with any SQLite client pointed at the Ratings service's
\code{ratings.db}, and send the page again:

\begin{minted}{sql}
ALTER TABLE Ratings RENAME TO RatingsMissing;
\end{minted}

It answers 200, every title and speaker
present and both rating fields null on all four sessions, and one error. Here
it is with the \code{data} half and six of the eight inner errors taken out,
because all eight say the same thing:

\begin{minted}[fontsize=\footnotesize,breakanywhere]{json}
{"errors":[{"message":"Failed to fetch from Subgraph 'ratings' at Path 'sessions.nodes'.","extensions":{"errors":[{"message":"Unexpected Execution Error","path":["sessions","nodes","averageScore"],"extensions":{"code":"DOWNSTREAM_SERVICE_ERROR"}},{"message":"Unexpected Execution Error","path":["sessions","nodes","ratingCount"],"extensions":{"code":"DOWNSTREAM_SERVICE_ERROR"}}],"serviceName":"ratings","statusCode":200}}]}
\end{minted}

\index{Unexpected Execution Error@\code{Unexpected Execution Error}!says nothing about the cause}%
Two things in it matter. \code{serviceName} is the router routing the failure
to its owner, in an extension a program can read, so the question
section~\ref{sec:ch18-the-page-that-got-slower} could not answer for a slow
page is answered for a failing one before anyone opens a log. And the eight
inner errors are Hot Chocolate's, and every one of them is
\code{Unexpected Execution Error}, two fields on four sessions, with nothing
about a table. The Ratings team now knows it is theirs and nothing else. The
cause is in the Ratings service's console, if it is anywhere, and the
response's \code{X-Wg-Trace-Id} is how they find the right part of it. Put
the table back the same way, \code{ALTER TABLE RatingsMissing RENAME TO
Ratings;}, before sending anything else.

\subsection{What the Owner Finds}

\index{Entity Framework Core!logs its own failures}%
For this failure the cause is there, and not because of anything in this
book. EF Core logs a command it failed to execute at error level, with the
\code{SqliteException} and its message, \code{no such table: Ratings}. With
section~\ref{sec:ch18-writing-it-down}'s scopes printed, those entries carry the
request's \code{TraceId}, and the statement marker before them carries the
same id, so the search the Speakers team ran in that section finds this
failure too.

\index{Hot Chocolate!writes a resolver's exception nowhere}%
That is luck about where the exception was thrown. Hot Chocolate turns an
exception inside a resolver into \code{Unexpected Execution Error} for the
caller and writes the exception itself nowhere, neither to the console nor
through the logger at the levels this book's services run at, with scopes on
or off. The database failure
above is found only because EF Core wrote its own entry before Hot Chocolate
swallowed the exception. Throw from a resolver that never touches the
database, and the router's \code{serviceName} sends the Ratings team to a
console holding not one line about the request.
Chapter~\ref{ch:second-third-subgraph} met this silence already, when the
Ratings service without its node id serializer answered
\code{Unexpected Execution Error} with nothing in any log. The trace id cannot
find an entry that was never written.

\subsection{A Listener That Writes It Down}

\index{diagnostic event listener!logging resolver errors|(}%
Hot Chocolate reports what happens during execution to \emph{diagnostic event
listeners}, classes that override the events they care about, and a resolver
error is one of those events, delivered with the \code{IError} the caller will
see and the exception behind it. So the fix is a listener that logs each one
through the service's logger, where the scope adds the trace id. Here is
\code{ErrorLog.cs} for the Ratings service, complete:

\begin{minted}{csharp}
using HotChocolate.Execution.Instrumentation;
using HotChocolate.Resolvers;

namespace Ratings;

/// <summary>
/// Writes every resolver error to this service's own log, with the
/// exception behind it. Hot Chocolate answers the caller with
/// Unexpected Execution Error and writes the exception nowhere, so
/// without this the one service that knows what went wrong keeps no
/// record of it. The log scope carries the trace id the router sent,
/// which is what lets a trace id from a bug report find this line.
/// </summary>
public sealed class ErrorLog(ILogger<ErrorLog> logger)
    : ExecutionDiagnosticEventListener
{
    public override void ResolverError(
        IMiddlewareContext context,
        IError error)
        => logger.LogError(
            error.Exception,
            "{Path}: {Message}",
            error.Path,
            error.Message);
}
\end{minted}

\index{Program.cs@\code{Program.cs}!registering ErrorLog}%
It is registered on the GraphQL builder, and the registration has a trap of
its own. Hot Chocolate builds listeners from the schema's own service
provider, and the error it gives shows that provider does not hold the host's
loggers. Registering the listener on its own compiles, and then
the service refuses to start, reporting that it cannot
resolve \code{ILogger<ErrorLog>} while activating \code{ErrorLog}.
Registering the logger with \code{AddApplicationService} first is what works. Here is the Ratings
service's \code{Program.cs}, complete, with the two registrations and the
comment that explains them marked:

\begin{minted}[highlightlines={47-53}]{csharp}
using ChilliCream.Nitro.App;
using HotChocolate.AspNetCore;
using Microsoft.EntityFrameworkCore;
using Ratings;

var builder = WebApplication.CreateBuilder(args);

// EF Core reports every command through the logging pipeline as well,
// at Information and with a duration attached. StatementLog below is
// this service's instrument, so silence the other one rather than have
// every statement printed twice.
builder.Logging.AddFilter(
    "Microsoft.EntityFrameworkCore.Database.Command",
    LogLevel.Warning);

builder.Services.AddDbContextFactory<RatingContext>(options =>
{
    options.UseSqlite("Data Source=ratings.db");

    if (builder.Environment.IsDevelopment())
    {
        options.AddInterceptors(new StatementLog());
    }
});

// Two lines here differ from the other two services, and both follow
// from the same fact: this service contributes fields to a type it
// does not own and owns no node of its own.
//
// There is no AddGlobalObjectIdentification, because there cannot be.
// That call refuses to build a schema with no node type in it: "There
// is no object type implementing interface `Node`." So this service
// has no node or nodes field, and needs no interceptor to declare
// them shareable.
//
// But the keys it is handed across the seam are still global node ids,
// and the serializer that reads one is normally registered by the call
// it cannot make. AddDefaultNodeIdSerializer registers it on its own.
// Without this line the reference resolver is injected a null and the
// entity route answers Unexpected Execution Error.
builder
    .AddGraphQL()
    .AddRatingsTypes()
    .RegisterDbContextFactory<RatingContext>()
    .AddDefaultNodeIdSerializer()
    .AddApolloFederation()
    // A resolver that throws is answered as Unexpected Execution Error
    // and Hot Chocolate writes the exception nowhere, so ErrorLog does.
    // Listeners are built from the schema's own services, which hold
    // none of the host's loggers: without the first line the service
    // refuses to start, unable to resolve ILogger<ErrorLog>.
    .AddApplicationService<ILogger<ErrorLog>>()
    .AddDiagnosticEventListener<ErrorLog>()
    .ModifyCostOptions(options => options.ApplyCostDefaults = false);

var app = builder.Build();

using (var scope = app.Services.CreateScope())
{
    var factory = scope.ServiceProvider
        .GetRequiredService<IDbContextFactory<RatingContext>>();
    using var db = factory.CreateDbContext();
    db.Database.EnsureCreated();
}

app.MapGraphQL()
    .WithOptions(options =>
    {
        // Serve the Nitro build that shipped in the package. The
        // default, ServeMode.Latest, downloads the current one
        // from ChilliCream's CDN instead.
        options.Tool.ServeMode = ServeMode.Embedded;
    });

app.RunWithGraphQLCommands(args);
\end{minted}

\index{ErrorLog@\code{ErrorLog}!what it writes}%
Rebuild and restart the Ratings service, rename the table again, send the
page, and put the table back afterwards. The Ratings console now holds one
\code{ErrorLog} entry per inner error, eight, each opening like this, with the
stack trace that follows cut:

\begin{minted}[fontsize=\footnotesize,breakanywhere]{text}
fail: Ratings.ErrorLog[0]
      => SpanId:161a2ce1c7f840fc, TraceId:262e3e16a3ca9dd8da4553a6834b0426, ParentId:d28132b2876aef8a => ConnectionId:0HNOI8R7TFM49 => RequestPath:/graphql RequestId:0HNOI8R7TFM49:00000004
      _entities[0].ratingCount: Unexpected Execution Error
      Microsoft.Data.Sqlite.SqliteException (0x80004005): SQLite Error 1: 'no such table: Ratings'.
\end{minted}

\index{trace id!joins a client's error to the exception}%
The second line is the scope, and \code{TraceId} in it is the value the client
received in \code{X-Wg-Trace-Id}. The third is the path inside the Ratings
service's own \code{\_entities} answer, which is not the path the client saw,
and the fourth is the exception the client was never shown. The listener logs
the database's failure a second time here, which is harmless. What it exists
for is the resolver that throws something else, where its entry is the only
one there is.
\index{diagnostic event listener!logging resolver errors|)}%

\subsection{The Same Lines in the Other Three Services}

\index{Program.cs@\code{Program.cs}!the same registration in four services}%
Every service in the graph gets \code{ErrorLog.cs}, identical but for its
namespace, and the same two registrations immediately before the call that
turns off cost defaults. A service that leaves them out is a service whose
failures leave no trace in its own log, however well the others are set up.
Here is the Sessions service's \code{Program.cs}, complete:

\begin{minted}[highlightlines={35-41}]{csharp}
using ChilliCream.Nitro.App;
using HotChocolate.AspNetCore;
using Microsoft.EntityFrameworkCore;
using Sessions;

var builder = WebApplication.CreateBuilder(args);

// EF Core reports every command through the logging pipeline as well,
// at Information and with a duration attached. StatementLog below is
// this service's instrument, so silence the other one rather than have
// every statement printed twice.
builder.Logging.AddFilter(
    "Microsoft.EntityFrameworkCore.Database.Command",
    LogLevel.Warning);

builder.Services.AddDbContextFactory<ConferenceContext>(options =>
{
    options.UseSqlite("Data Source=conference.db");

    if (builder.Environment.IsDevelopment())
    {
        options.AddInterceptors(new StatementLog());
    }
});

builder
    .AddGraphQL()
    .AddSessionsTypes()
    .RegisterDbContextFactory<ConferenceContext>()
    .AddQueryContext()
    .AddPagingArguments()
    .AddGlobalObjectIdentification()
    .AddMutationConventions()
    .AddApolloFederation()
    // A resolver that throws is answered as Unexpected Execution Error
    // and Hot Chocolate writes the exception nowhere, so ErrorLog does.
    // Listeners are built from the schema's own services, which hold
    // none of the host's loggers: without the first line the service
    // refuses to start, unable to resolve ILogger<ErrorLog>.
    .AddApplicationService<ILogger<ErrorLog>>()
    .AddDiagnosticEventListener<ErrorLog>()
    .ModifyCostOptions(options => options.ApplyCostDefaults = false)
    .TryAddTypeInterceptor<ShareNodeFields>()
    .TryAddTypeInterceptor<NullableNodesField>();

var app = builder.Build();

using (var scope = app.Services.CreateScope())
{
    var factory = scope.ServiceProvider
        .GetRequiredService<IDbContextFactory<ConferenceContext>>();
    using var db = factory.CreateDbContext();
    db.Database.EnsureCreated();
}

app.MapGraphQL()
    .WithOptions(options =>
    {
        // Serve the Nitro build that shipped in the package. The
        // default, ServeMode.Latest, downloads the current one
        // from ChilliCream's CDN instead.
        options.Tool.ServeMode = ServeMode.Embedded;
    });

app.RunWithGraphQLCommands(args);
\end{minted}

The Speakers service's, which carries chapter~\ref{ch:entities-authorization}'s
authentication as well:

\begin{minted}[highlightlines={63-69}]{csharp}
using System.Text;
using ChilliCream.Nitro.App;
using HotChocolate.AspNetCore;
using Microsoft.EntityFrameworkCore;
using Microsoft.IdentityModel.Tokens;
using Speakers;

var builder = WebApplication.CreateBuilder(args);

// EF Core reports every command through the logging pipeline as well,
// at Information and with a duration attached. StatementLog below is
// this service's instrument, so silence the other one rather than have
// every statement printed twice.
builder.Logging.AddFilter(
    "Microsoft.EntityFrameworkCore.Database.Command",
    LogLevel.Warning);

builder.Services.AddDbContextFactory<SpeakerContext>(options =>
{
    options.UseSqlite("Data Source=speakers.db");

    if (builder.Environment.IsDevelopment())
    {
        options.AddInterceptors(new StatementLog());
    }
});

// A shared secret rather than a JWKS endpoint, because this book
// stands up no identity provider and a reader has to be able to mint
// a token and repeat every request in the chapter. The three
// Validate flags are off for the same reason and are the three a
// deployment turns on first. What is not negotiable is that the
// router validates the same token against the same key: a subgraph
// that trusts a header the router set, without checking it, has
// moved the guard back out of the service again.
builder.Services
    .AddAuthentication()
    .AddJwtBearer(options =>
    {
        options.TokenValidationParameters = new()
        {
            ValidateIssuer = false,
            ValidateAudience = false,
            ValidateLifetime = false,
            IssuerSigningKey = new SymmetricSecurityKey(
                Encoding.UTF8.GetBytes(SigningKey.Value))
        };
    });

builder.Services.AddAuthorization();

builder
    .AddGraphQL()
    // Without this the [Authorize] attributes compile, reach the
    // schema as @authorize, and are never executed.
    .AddAuthorization()
    .AddSpeakersTypes()
    .RegisterDbContextFactory<SpeakerContext>()
    .AddQueryContext()
    .AddPagingArguments()
    .AddGlobalObjectIdentification()
    .AddApolloFederation()
    // A resolver that throws is answered as Unexpected Execution Error
    // and Hot Chocolate writes the exception nowhere, so ErrorLog does.
    // Listeners are built from the schema's own services, which hold
    // none of the host's loggers: without the first line the service
    // refuses to start, unable to resolve ILogger<ErrorLog>.
    .AddApplicationService<ILogger<ErrorLog>>()
    .AddDiagnosticEventListener<ErrorLog>()
    .ModifyCostOptions(options => options.ApplyCostDefaults = false)
    .TryAddTypeInterceptor<ShareNodeFields>()
    .TryAddTypeInterceptor<NullableNodesField>();

var app = builder.Build();

using (var scope = app.Services.CreateScope())
{
    var factory = scope.ServiceProvider
        .GetRequiredService<IDbContextFactory<SpeakerContext>>();
    using var db = factory.CreateDbContext();
    db.Database.EnsureCreated();
}

app.UseAuthentication();
app.UseAuthorization();

app.MapGraphQL()
    .WithOptions(options =>
    {
        // Serve the Nitro build that shipped in the package. The
        // default, ServeMode.Latest, downloads the current one
        // from ChilliCream's CDN instead.
        options.Tool.ServeMode = ServeMode.Embedded;
    });

app.RunWithGraphQLCommands(args);
\end{minted}

And the Search service's:

\begin{minted}[highlightlines={28-34}]{csharp}
using ChilliCream.Nitro.App;
using HotChocolate.AspNetCore;
using Microsoft.EntityFrameworkCore;
using Search;

var builder = WebApplication.CreateBuilder(args);

builder.Logging.AddFilter(
    "Microsoft.EntityFrameworkCore.Database.Command",
    LogLevel.Warning);

builder.Services.AddDbContextFactory<SearchContext>(options =>
{
    options.UseSqlite("Data Source=search.db");
    if (builder.Environment.IsDevelopment())
    {
        options.AddInterceptors(new StatementLog());
    }
});

builder
    .AddGraphQL()
    .AddSearchTypes()
    .RegisterDbContextFactory<SearchContext>()
    .AddPagingArguments()
    .AddDefaultNodeIdSerializer()
    .AddApolloFederation()
    // A resolver that throws is answered as Unexpected Execution Error
    // and Hot Chocolate writes the exception nowhere, so ErrorLog does.
    // Listeners are built from the schema's own services, which hold
    // none of the host's loggers: without the first line the service
    // refuses to start, unable to resolve ILogger<ErrorLog>.
    .AddApplicationService<ILogger<ErrorLog>>()
    .AddDiagnosticEventListener<ErrorLog>()
    .ModifyCostOptions(options => options.ApplyCostDefaults = false);

var app = builder.Build();

using (var scope = app.Services.CreateScope())
{
    var factory = scope.ServiceProvider
        .GetRequiredService<IDbContextFactory<SearchContext>>();
    using var db = factory.CreateDbContext();
    db.Database.EnsureCreated();
}

app.MapGraphQL().WithOptions(options =>
{
    options.Tool.ServeMode = ServeMode.Embedded;
});

app.RunWithGraphQLCommands(args);
\end{minted}

\index{blame routing!the router routes, the service explains}%
The router routes a failure to a service by name, and a slow request to
nobody. What turns the id into lines an owner can find is the router handing
it out, and each service owner making section~\ref{sec:ch18-writing-it-down}'s
two service edits and adding the listener above.
\index{blame routing!a failure, routed|)}%
